%% Reference list reformatted to Journal of Computational Finance (Risk Journals)
%% house style: alphabetical by author; "Last, F., and Last, F. (Year). Title.
%% Journal vol(issue), pages." Include DOIs where available. In-text citations are
%% author-date, e.g. "(Hagan and West 2006)" or "Hagan and West (2006)".
%%
%% Drop this list into the JCF manuscript (replacing the numbered \thebibliography),
%% and switch in-text \cite{...} to author-date form. See JCF_submission.md.

\begin{thebibliography}{}

\bibitem{Ametrano} Ametrano, F. M., and Henrard, M. (2014). Everything you always
wanted to know about multiple interest rate curve bootstrapping but were afraid to ask.
Working Paper. URL: https://ssrn.com/abstract=2219548.

\bibitem{Anderson} Andersen, L. (2007). Discount curve construction with tension
splines. Review of Derivatives Research 10(3), 227--267.
https://doi.org/10.1007/s11147-008-9021-2.

\bibitem{Bianchetti} Bianchetti, M. (2010). Two curves, one price. Risk 23(8), 66--72.

\bibitem{Cline} Cline, A. K. (1974). Scalar- and planar-valued curve fitting using
splines under tension. Communications of the ACM 17(4), 218--223.
https://doi.org/10.1145/360924.360971.

\bibitem{Costantini} Costantini, P., Kvasov, B. I., and Manni, C. (1999). On discrete
hyperbolic tension splines. Advances in Computational Mathematics 11(4), 331--354.
https://doi.org/10.1023/A:1018988312596.

\bibitem{Fritsch} Fritsch, F. N., and Carlson, R. E. (1980). Monotone piecewise cubic
interpolation. SIAM Journal on Numerical Analysis 17(2), 238--246.
https://doi.org/10.1137/0717021.

\bibitem{Hagan} Hagan, P. S., and West, G. (2006). Interpolation methods for curve
construction. Applied Mathematical Finance 13(2), 89--129.
https://doi.org/10.1080/13504860500396032.

\bibitem{H15} Board of Governors of the Federal Reserve System (2026). Selected
Interest Rates (H.15). Daily release, May 22.
URL: https://www.federalreserve.gov/releases/h15/.

\bibitem{Henrard} Henrard, M. (2014). Interest Rate Modelling in the Multi-Curve
Framework: Foundations, Evolution and Implementation. Applied Quantitative Finance.
Palgrave Macmillan. https://doi.org/10.1057/9781137374660.

\bibitem{Kamada} Kamada, M., and Enkhbat, R. (2009). Spline interpolation in piecewise
constant tension. In Proceedings of the 8th International Conference on Sampling Theory
and Applications (SampTA'09), Marseille.

\bibitem{Kincaid} Kincaid, D., and Cheney, W. (2002). Numerical Analysis: Mathematics
of Scientific Computing, 3rd edn. American Mathematical Society.

\bibitem{McCulloch} McCulloch, J. H. (1971). Measuring the term structure of interest
rates. Journal of Business 44(1), 19--31. https://doi.org/10.1086/295329.

\bibitem{Renka} Renka, R. J. (1987). Interpolatory tension splines with automatic
selection of tension factors. SIAM Journal on Scientific and Statistical Computing
8(3), 393--415. https://doi.org/10.1137/0908041.

\bibitem{Ron} Ron, U. (2000). A practical guide to swap curve construction. Working
Paper 2000-17, Bank of Canada.

\bibitem{Schweikert} Schweikert, D. G. (1966). An interpolation curve using a spline in
tension. Journal of Mathematics and Physics 45(1--4), 312--317.
https://doi.org/10.1002/sapm1966451312.

\bibitem{Tanggaard} Tanggaard, C. (1997). Nonparametric smoothing of yield curves.
Review of Quantitative Finance and Accounting 9(3), 251--267.
https://doi.org/10.1023/A:1008286426806.

\bibitem{TreasuryMethod} US Department of the Treasury (2021). Treasury par yield curve
methodology. Technical Note (monotone convex spline, effective December 6).
URL: https://home.treasury.gov/.

\end{thebibliography}
